\chapter{2011年安徽卷（文）}

\section{单选题}

\begin{problem}
设 \(\i\) 是虚数单位，复数 \(\frac{1 + a\i}{2 - \i}\) 为纯虚数，则实数 \(a\) 为（\quad）

\choices
    {\(2\)}
    {\(-2\)}
    {\(-\frac{1}{2}\)}
    {\(\frac{1}{2}\)}
\end{problem}

\begin{answer}
A
\end{answer}

\begin{problem}
集合 \(U = \{1, 2, 3, 4, 5, 6\}\)，\(S = \{1, 4, 5\}\)，\(T = \{2, 3, 4\}\)，则 \(S \cap \left(\complement_{U} T\right)\) 等于（\quad）

\choices
    {\(\{1, 4, 5, 6\}\)}
    {\(\{1, 5\}\)}
    {\(\{4\}\)}
    {\(\{1, 2, 3, 4, 5\}\)}
\end{problem}

\begin{answer}
B
\end{answer}

\begin{problem}
双曲线 \(2x^{2} - y^{2} = 8\) 的实轴长是（\quad）

\choices
    {\(2\)}
    {\(2\sqrt{2}\)}
    {\(4\)}
    {\(4\sqrt{2}\)}
\end{problem}

\begin{answer}
C
\end{answer}

\begin{problem}
若直线 \(3x + y + a = 0\) 过圆 \(x^{2} + y^{2} + 2x - 4y = 0\) 的圆心，则 \(a\) 的值为（\quad）

\choices
    {\(-1\)}
    {\(1\)}
    {\(3\)}
    {\(-3\)}
\end{problem}

\begin{answer}
B
\end{answer}

\begin{problem}
若点 \((a, b)\) 在 \(y = \lg x\) 图象上，\(a \neq 1\)，则下列点中在此图象上的是（\quad）

\choices
    {\(\left(\frac{1}{a}, b\right)\)}
    {\((10a, 1 - b)\)}
    {\(\left(\frac{10}{a}, b + 1\right)\)}
    {\((a^{2}, 2b)\)}
\end{problem}

\begin{answer}
D
\end{answer}

\begin{problem}
设变量 \(x\)，\(y\) 满足 \(\begin{cases}
x + y \leqslant 1,\\
x - y \leqslant 1,\\
x \geqslant 0,
\end{cases}\)
则 \(x + 2y\) 的最大值和最小值分别为（\quad）

\choices
    {\(1\)，\(-1\)}
    {\(2\)，\(-2\)}
    {\(1\)，\(-2\)}
    {\(2\)，\(-1\)}
\end{problem}

\begin{answer}
B
\end{answer}

\begin{problem}
若数列 \(\{a_{n}\}\) 的通项公式是 \(a_{n} = (-1)^{n}(3n - 2)\)，则 \(a_{1} + a_{2} + \cdots + a_{10} =\)（\quad）

\choices
    {\(15\)}
    {\(12\)}
    {\(-12\)}
    {\(-15\)}
\end{problem}

\begin{answer}
A
\end{answer}

\begin{problem}
一个空间几何体的三视图如图所示，则该几何体的表面积为（\quad）

\bitmapfigure[max width=0.32\linewidth]{img/2011/anhui_liberal/q8_fig1.png}

\choices
    {\(48\)}
    {\(32 + 8\sqrt{17}\)}
    {\(48 + 8\sqrt{17}\)}
    {\(80\)}
\end{problem}

\begin{answer}
C
\end{answer}

\begin{problem}
从正六边形的 6 个顶点中随机选择 4 个顶点，则以它们作为顶点的四边形是矩形的概率等于（\quad）

\choices
    {\(\frac{1}{10}\)}
    {\(\frac{1}{8}\)}
    {\(\frac{1}{6}\)}
    {\(\frac{1}{5}\)}
\end{problem}

\begin{answer}
D
\end{answer}

\begin{problem}
函数 \(f(x) = ax^{n}(1 - x)^{2}\) 在区间 \([0, 1]\) 上的图象如图所示，则 \(n\) 可能是（\quad）

\bitmapfigure[max width=0.28\linewidth]{img/2011/anhui_liberal/q10_fig1.png}

\choices
    {\(1\)}
    {\(2\)}
    {\(3\)}
    {\(4\)}
\end{problem}

\begin{answer}
A
\end{answer}

\section{填空题}

\begin{problem}
设 \(f(x)\) 是定义在 \(\R\) 上的奇函数，当 \(x \leqslant 0\) 时，\(f(x) = 2x^{2} - x\)，则 \(f(1) =\fillinblank{}\).
\end{problem}

\begin{answer}
\(-3\)
\end{answer}

\begin{problem}
如图所示，程序框图（算法流程图）的输出结果是\fillinblank{}.

\bitmapfigure[float=inline,max width=0.5\linewidth]{img/2011/anhui_liberal/q12_fig1.png}
\end{problem}

\begin{answer}
\(15\)
\end{answer}

\begin{problem}
函数 \(y = \frac{1}{\sqrt{6 - x - x^{2}}}\) 的定义域是\fillinblank{}.
\end{problem}

\begin{answer}
\(\left(-3,2\right)\)
\end{answer}

\begin{problem}
已知向量 \(\bs{a}\)，\(\bs{b}\) 满足 \((\bs{a} + 2\bs{b}) \cdot (\bs{a} - \bs{b}) = -6\)，且 \(|\bs{a}| = 1\)，\(|\bs{b}| = 2\)，则 \(\bs{a}\) 与 \(\bs{b}\) 的夹角为\fillinblank{}.
\end{problem}

\begin{answer}
\(\frac{\pi}{3}\)
\end{answer}

\begin{problem}
设 \(f(x) = a\sin 2x + b\cos 2x\)，其中 \(a\)，\(b \in \R\)，\(ab \neq 0\)，若 \(f(x) \leqslant \left|f\left(\frac{\pi}{6}\right)\right|\) 对一切 \(x \in \R\) 恒成立，则
\begin{circlelist}
\item \(f\left(\frac{11\pi}{12}\right) = 0\)；
\item \(f\left(\frac{7\pi}{10}\right) < \left|f\left(\frac{\pi}{5}\right)\right|\)；
\item \(f(x)\) 既不是奇函数，也不是偶函数；
\item \(f(x)\) 的单调递增区间是 \(\left[k\pi + \frac{\pi}{6}, k\pi + \frac{2\pi}{3}\right]\)（\(k \in \Z\)）；
\item 存在经过点 \((a, b)\) 的直线与函数 \(f(x)\) 的图象不相交.
\end{circlelist}
以上结论正确的是\fillinblank{}.（写出所有正确结论的编号）
\end{problem}

\begin{answer}
\(1\)，\(3\)
\end{answer}

\section{解答题}

\begin{problem}
在 \(\triangle ABC\) 中，\(a\)，\(b\)，\(c\) 分别为内角 \(A\)，\(B\)，\(C\) 所对的边长，\(a = \sqrt{3}\)，\(b = \sqrt{2}\)，\(1 + 2\cos(B + C) = 0\)，求边 \(BC\) 上的高.
\end{problem}

\begin{solution}
由 \(B+C\in(0,\pi)\) 且 \(1+2\cos(B+C)=0\)，得 \(\cos(B+C)=-\frac12\)，所以 \(B+C=\frac{2\pi}{3}\)，从而 \(A=\frac\pi3\).

由正弦定理，\(\sin B=\frac{b\sin A}{a}=\frac{\sqrt2\cdot\frac{\sqrt3}{2}}{\sqrt3}=\frac{\sqrt2}{2}\).
又因为 \(0<B<B+C=\frac{2\pi}{3}\)，故 \(B=\frac\pi4\)，进而 \(C=\frac{5\pi}{12}\).

再次由正弦定理，\(c=\frac{a\sin C}{\sin A}=\frac{\sqrt3\sin\frac{5\pi}{12}}{\sin\frac\pi3}=2\sin\frac{5\pi}{12}=\frac{\sqrt6+\sqrt2}{2}\). 设 \(A\) 到 \(BC\) 的高为 \(h\)，则在直角三角形中 \(h=c\sin B\)，所以 \(h=\frac{\sqrt6+\sqrt2}{2}\cdot\frac{\sqrt2}{2}=\frac{\sqrt3+1}{2}\).
故边 \(BC\) 上的高为 \(\frac{\sqrt3+1}{2}\).
\end{solution}

\begin{problem}
设直线 \(l_{1} : y = k_{1}x + 1\)，\(l_{2} : y = k_{2}x - 1\)，其中实数 \(k_{1}\)，\(k_{2}\) 满足 \(k_{1}k_{2} + 2 = 0\).
\begin{enumerate}
\item 证明 \(l_{1}\) 与 \(l_{2}\) 相交；
\item 证明 \(l_{1}\) 与 \(l_{2}\) 的交点在椭圆 \(2x^{2} + y^{2} = 1\) 上.
\end{enumerate}
\end{problem}

\begin{solution}
\begin{enumerate}
\item 由 \(k_{1}k_{2}+2=0\) 得 \(k_{1}k_{2}=-2\)，所以 \(k_{1}\ne k_{2}\). 联立两直线方程，得 \(k_{1}x+1=k_{2}x-1\)，即 \((k_{1}-k_{2})x=-2\).
因为 \(k_{1}-k_{2}\ne0\)，上式有唯一实数解 \(x\)，再由任一直线方程唯一确定 \(y\)，所以 \(l_{1}\) 与 \(l_{2}\) 相交.

\item 设交点为 \(P(x,y)\). 由 \(P\in l_{1}\) 和 \(P\in l_{2}\)，分别有 \(y-1=k_{1}x\)，\(y+1=k_{2}x\).
两式相乘，得
\[
y^{2}-1=(y-1)(y+1)=k_{1}k_{2}x^{2}=-2x^{2},
\]
即 \(2x^{2}+y^{2}=1\). 因此交点 \(P\) 在椭圆 \(2x^{2}+y^{2}=1\) 上.
\end{enumerate}
\end{solution}

\begin{problem}
设 \(f(x) = \frac{\e^{x}}{1 + ax^{2}}\)，其中 \(a\) 为正实数.
\begin{enumerate}
\item 当 \(a = \frac{4}{3}\) 时，求 \(f(x)\) 的极值点；
\item 若 \(f(x)\) 为 \(\R\) 上的单调函数，求 \(a\) 的取值范围.
\end{enumerate}
\end{problem}

\begin{solution}
\begin{enumerate}
\item 求导得
\[
f'(x)=\frac{\e^{x}(1+ax^{2})-\e^{x}\cdot2ax}{(1+ax^{2})^{2}}
=\frac{\e^{x}\left[a(x-1)^{2}+1-a\right]}{(1+ax^{2})^{2}}.
\]
当 \(a=\frac43\) 时，
\[
a(x-1)^{2}+1-a=\frac43(x-1)^{2}-\frac13=\frac13(2x-1)(2x-3).
\]
由于 \(\e^{x}>0\) 且 \((1+ax^{2})^{2}>0\)，所以 \(f'(x)\) 在区间 \(\left(-\infty,\frac12\right)\) 和 \(\left(\frac32,+\infty\right)\) 上为正，在 \(\left(\frac12,\frac32\right)\) 上为负. 因此 \(f(x)\) 在 \(x=\frac12\) 处取得极大值，在 \(x=\frac32\) 处取得极小值，且
\(f\left(\frac12\right)=\frac{\e^{1/2}}{1+\frac43\cdot\frac14}=\frac{3\sqrt{\e}}4\)，\(f\left(\frac32\right)=\frac{\e^{3/2}}{1+\frac43\cdot\frac94}=\frac{\e^{3/2}}4\).
故极大值点为 \(\left(\frac12,\frac{3\sqrt{\e}}4\right)\)，极小值点为 \(\left(\frac32,\frac{\e^{3/2}}4\right)\).

\item 因为前式中除去正因子 \(\frac{\e^{x}}{(1+ax^{2})^{2}}\) 后，导数的符号由 \(q(x)=a(x-1)^{2}+1-a\) 决定. 当 \(0<a\leqslant1\) 时，对任意 \(x\in\R\) 都有 \(q(x)\geqslant1-a\geqslant0\)，所以 \(f'(x)\geqslant0\)，函数在 \(\R\) 上单调递增.

当 \(a>1\) 时，\(q(1)=1-a<0\)，而当 \(\lvert x-1\rvert\) 足够大时 \(q(x)>0\)，故 \(f'(x)\) 先为正、后为负、再为正，函数不可能在 \(\R\) 上单调. 因而所求 \(a\) 的取值范围为 \(0<a\leqslant1\).
\end{enumerate}
\end{solution}

\begin{problem}
如图，\(ABEDFC\) 为多面体，平面 \(ABED\) 与平面 \(ACFD\) 垂直，点 \(O\) 在线段 \(AD\) 上，\(OA = 1\)，\(OD = 2\)，\(\triangle OAB\)，\(\triangle OAC\)，\(\triangle ODE\)，\(\triangle ODF\) 都是正三角形.

\begin{enumerate}
\item 证明直线 \(BC \parallel EF\)；
\item 求棱锥 \(F - OBED\) 的体积.
\end{enumerate}

\bitmapfigure[max width=0.34\linewidth]{img/2011/anhui_liberal/q19_fig1.png}
\end{problem}

\begin{solution}
取 \(O\) 为原点，令 \(AD\) 为 \(x\) 轴，平面 \(ABED\) 为 \(xOy\) 平面，平面 \(ACFD\) 为 \(xOz\) 平面. 由 \(OA=1\)，\(OD=2\)，可取 \(A=(-1,0,0)\)，\(D=(2,0,0)\). 按图中各正三角形所在的一侧取点，则 \(B=\left(-\frac12,\frac{\sqrt3}{2},0\right)\)，\(E=(1,\sqrt3,0)\)，\(C=\left(-\frac12,0,\frac{\sqrt3}{2}\right)\)，\(F=(1,0,\sqrt3)\).
这些坐标分别满足 \(OA=OB=AB=1\)、\(OD=OE=DE=2\) 以及相应的正三角形条件.

\begin{enumerate}
\item 有
\[
\overrightarrow{BC}=C-B=\left(0,-\frac{\sqrt3}{2},\frac{\sqrt3}{2}\right),
\]
\(\overrightarrow{EF}=F-E=(0,-\sqrt3,\sqrt3)=2\overrightarrow{BC}\).
所以 \(BC\parallel EF\).

\item 在平面 \(ABED\) 内，
\(\overrightarrow{OB}=\left(-\frac12,\frac{\sqrt3}{2},0\right)\)，\(\overrightarrow{DE}=(-1,\sqrt3,0)=2\overrightarrow{OB}\). 故四边形 \(OBED\) 是以 \(OB\) 和 \(DE\) 为底的梯形，且 \(|OB|=1\)，\(|DE|=2\). 两底所在平行直线之间的距离为 \(OD\sin\angle BOD=2\sin\frac{2\pi}{3}=\sqrt3\). 因此 \(S_{OBED}=\frac{1+2}{2}\cdot\sqrt3=\frac{3\sqrt3}{2}\).
点 \(F\) 到平面 \(ABED\) 的距离为其 \(z\) 坐标的绝对值，即 \(\sqrt3\). 所以棱锥 \(F-OBED\) 的体积为
\[
V=\frac13S_{OBED}\cdot\sqrt3
=\frac13\cdot\frac{3\sqrt3}{2}\cdot\sqrt3=\frac32.
\]
\end{enumerate}
\end{solution}

\begin{problem}
某地最近十年粮食需求量逐年上升，下表是部分统计数据：

\begin{center}
\begin{tabular}{c|ccccc}
年份 & \(2002\) & \(2004\) & \(2006\) & \(2008\) & \(2010\) \\
\hline
需求量（\(\text{万吨}\)） & \(236\) & \(246\) & \(257\) & \(276\) & \(286\)
\end{tabular}
\end{center}

\begin{enumerate}
\item 利用所给数据求年需求量与年份之间的回归直线方程 \(\hat{y} = bx + a\)；
\item 利用（1）中所求出的直线方程预测该地 \(2012\) 年的粮食需求量.
\end{enumerate}
\end{problem}

\begin{solution}
\begin{enumerate}
\item 记年份为 \(x\)，需求量为 \(y\). 由所给数据，
\[
\bar{x}=\frac{2002+2004+2006+2008+2010}{5}=2006.
\]
\[
\bar{y}=\frac{236+246+257+276+286}{5}=260.2.
\]
令 \(u=x-2006\)，则五个 \(u\) 值依次为 \(-4,-2,0,2,4\). 因而
\[
\sum u^{2}=(-4)^{2}+(-2)^{2}+0^{2}+2^{2}+4^{2}=40,
\]
\[
\begin{gathered}
\sum u(y-\bar y)=(-4)(236-260.2)+(-2)(246-260.2)\\
\quad+2(276-260.2)+4(286-260.2)=260.
\end{gathered}
\]
所以回归系数
\(b=\frac{\sum (x-\bar x)(y-\bar y)}{\sum (x-\bar x)^{2}}=\frac{260}{40}=6.5\)，\(a=\bar y-b\bar x=260.2-6.5\times2006=-12778.8\).
故回归直线方程为
\[
\hat y=6.5x-12778.8=6.5(x-2006)+260.2.
\]

\item 当 \(x=2012\) 时，
\[
\hat y=6.5(2012-2006)+260.2=6.5\times6+260.2=299.2.
\]
所以预测该地 \(2012\) 年的粮食需求量约为 \(299.2\text{ 万吨}\).
\end{enumerate}
\end{solution}

\begin{problem}
在数 \(1\) 和 \(100\) 之间插入 \(n\) 个实数，使得这 \(n + 2\) 个数构成递增的等比数列，将这 \(n + 2\) 个数的乘积记作 \(T_{n}\)，再令 \(a_{n} = \lg T_{n}\)，\(n \geqslant 1\).
\begin{enumerate}
\item 求数列 \(\{a_{n}\}\) 的通项公式；
\item 设 \(b_{n} = \tan a_{n} \cdot \tan a_{n+1}\)，求数列 \(\{b_{n}\}\) 的前 \(n\) 项和 \(S_{n}\).
\end{enumerate}
\end{problem}

\begin{solution}
\begin{enumerate}
\item 设这个递增等比数列的公比为 \(q\). 由于首项为 \(1\)，末项为 \(100\)，共有 \(n+2\) 项，所以
\[
q^{n+1}=100=10^{2}.
\]
又因数列递增，\(q>1\)，从而
\[
q=10^{\frac{2}{n+1}}.
\]
数列各项为 \(1,q,q^{2},\ldots,q^{n+1}\)，故其乘积为
\[
T_n=q^{0+1+\cdots+(n+1)}=q^{\frac{(n+1)(n+2)}{2}}
=\left(10^{\frac{2}{n+1}}\right)^{\frac{(n+1)(n+2)}{2}}=10^{n+2}.
\]
因此
\[
a_n=\lg T_n=\lg 10^{n+2}=n+2.
\]

\item 由上式，\(a_j=j+2\)，所以
\[
a_j-a_{j+1}=-1.
\]
利用正切差角公式，
\[
\tan(-1)=\tan(a_j-a_{j+1})
=\frac{\tan a_j-\tan a_{j+1}}{1+\tan a_j\tan a_{j+1}}.
\]
由此得
\[
\tan a_j\tan a_{j+1}
=\frac{\tan a_j-\tan a_{j+1}}{\tan(-1)}-1.
\]
于是
\[
S_n=\sum_{j=1}^{n}\tan a_j\tan a_{j+1}=\frac{1}{\tan(-1)}\sum_{j=1}^{n}(\tan a_j-\tan a_{j+1})-n=\frac{\tan a_1-\tan a_{n+1}}{\tan(-1)}-n.
\]
代入 \(a_1=3\)、\(a_{n+1}=n+3\) 以及 \(\tan(-1)=-\tan1\)，得
\[
S_n=\frac{\tan(n+3)-\tan3}{\tan1}-n.
\]
\end{enumerate}
\end{solution}
